%% kapa: exam=1 round=36 year=2025 period=2 form=A subject=appraisal_law subject_ko=감정평가관계법규 seq=38 number=38
%% source: Q-Net 감정평가사 기출문제 (cst00309) artlSeq=5231525 file=raw/1008/5231525/2025년 제36회 감정평가사 1차시험 2교시 문제지 원본.pdf
%% conversion: kordoc 4.24.51 -> tools/gen1.py (auto); review=proofread-sonnet; flags=none
\begin{question}{38}
도시 및 주거환경정비법령상 관리처분계획 등에 관한 설명으로 옳지 않은 것은? (단, 다른 조건은 고려하지 않음)
\begin{choices}
  \item 사업시행자는 분양신청을 받은 후 잔여분이 있는 경우에는 정관으로 정한 목적을 위하여 그 잔여분을 보류지로 정할 수 있다.
  \item 사업시행자는 정비사업의 시행으로 건설된 건축물을 인가받은 관리처분계획에 따라 토지등소유자에게 공급하여야 한다.
  \item 시장은 관리처분계획의 타당성 검증을 요청하는 경우에는 관리처분계획인가의 신청을 받은 날부터 60일 이내에 인가 여부를 결정하여 사업시행자에게 통지하여야 한다.
  \item 사업시행자는 분양신청을 하지 아니한 자와 손실보상에 관한 협의가 성립되지 아니하면 협의기간의 만료일 다음 날부터 60일 이내에 수용재결을 신청하거나 매도청구소송을 제기하여야 한다.
  \item 종전의 건축물의 전세권자는 사업시행자의 동의를 받지 않더라도 관리처분계획인가의 고시가 있은 때에는 이전고시가 있는 날까지 계속 사용하거나 수익할 수 있다.
\end{choices}
\end{question}
\answer{5}
